
\documentclass{article}
\usepackage{colortbl}
\usepackage{makecell}
\usepackage{multirow}
\usepackage{supertabular}

\begin{document}

\newcounter{utterance}

\centering \large Interaction Transcript for game `eqbench', experiment `eq\_bench', episode 93 with reflexion{-}llama{-}3.1{-}8b{-}instruct.
\vspace{24pt}

{ \footnotesize  \setcounter{utterance}{1}
\setlength{\tabcolsep}{0pt}
\begin{supertabular}{c@{$\;$}|p{.15\linewidth}@{}p{.15\linewidth}p{.15\linewidth}p{.15\linewidth}p{.15\linewidth}p{.15\linewidth}}
    \# & \multicolumn{2}{c}{Player} && \multicolumn{2}{c}{Game Master} \\
    \hline

    \theutterance \stepcounter{utterance}  
    & & & \multicolumn{4}{p{0.6\linewidth}}{
        \cellcolor[rgb]{0.9,0.9,0.9}{
            \makecell[{{p{\linewidth}}}]{
                \texttt{\tiny{[P1$\langle$GM]}}
                \texttt{You have attempted this same task before and did not succeed. These are the plans you wrote after those attempts.} \\
\texttt{Trial \#1: New plan: } \\
\\ 
\texttt{In my previous attempt, I failed to accurately predict Emma's emotional response. Upon re{-}reading the dialogue, I realized that I didn't consider the context of the conversation and the relationships involved. Emma is married to Steve, and Sean's confession of love puts her in a difficult position. } \\
\\ 
\texttt{To improve my prediction, I will focus on the following key points:} \\
\\ 
\texttt{1.  I will analyze the dialogue from Emma's perspective, considering her relationships and the potential consequences of her actions.} \\
\texttt{2.  I will pay attention to the emotional cues in the dialogue, such as Emma's hesitation and her attempts to clarify Sean's intentions.} \\
\texttt{3.  I will consider the power dynamics at play in the conversation, including Emma's marriage and Sean's confession of love.} \\
\\ 
\texttt{With these factors in mind, I will reassess the emotional response of Emma at the end of the dialogue.} \\
\\ 
\\ 
\\ 
\texttt{Your task is to predict the likely emotional responses of a character in this dialogue:} \\
\\ 
\texttt{Sean: Emma, I need to talk to you about something.} \\
\texttt{Emma: Of course, Sean. What's on your mind?} \\
\texttt{Sean: It's...it's about us. I mean, you and me.} \\
\texttt{Emma: Us? What do you mean?} \\
\texttt{Sean: I mean, I love you, Emma. I'm in love with you.} \\
\texttt{Emma: Sean...I...} \\
\texttt{Sean: Don't say anything yet. Just hear me out. I've been feeling this way for a while now. I can't ignore it anymore.} \\
\texttt{Emma: Sean, I'm married. You know that.} \\
\texttt{Sean: I do know that, Emma. But I can't help how I feel.} \\
\texttt{Emma: This is...this is a lot to take in.} \\
\texttt{Sean: I didn't mean to overwhelm you. I just...I needed to tell you.} \\
\texttt{Emma: Sean, I care about you, but not in the way you want me to.} \\
\texttt{Sean: So, that's it then? You're just going to dismiss my feelings?} \\
\texttt{Emma: No, Sean, I'm not dismissing your feelings. But I'm married. I love Steve. And I can't...I can't give you what you want.} \\
\texttt{Sean: I never asked you to leave Steve.} \\
\texttt{Emma: But you're asking me to betray him. To betray my marriage.} \\
\texttt{Sean: That's not what I meant.} \\
\texttt{Emma: But that's what it would be, Sean. And I can't do that.} \\
\texttt{Sean: So, what happens now?} \\
\texttt{Emma: I...I don't know, Sean. I really don't.} \\
\texttt{Sean: Maybe it's best if I find a new trainer.} \\
\texttt{Emma: Sean, don't be rash.} \\
\texttt{Sean: I don't think I can be around you right now, Emma. Not without...not without feeling this way.} \\
\texttt{{[}End dialogue{]}} \\
\\ 
\texttt{At the end of this dialogue, Emma would feel...} \\
\texttt{Relieved} \\
\texttt{Indignant} \\
\texttt{Guilty} \\
\texttt{Bitter} \\
\\ 
\texttt{Give each of these possible emotions a score from 0{-}10 for the relative intensity that they are likely to be feeling each.} \\
\\ 
\texttt{You must output in the following format, including headings (of course, you should give your own scores), with no additional commentary:} \\
\\ 
\texttt{Relieved: <score>} \\
\texttt{Indignant: <score>} \\
\texttt{Guilty: <score>} \\
\texttt{Bitter: <score>} \\
\\ 
\\ 
\texttt{{[}End of answer{]}} \\
\\ 
\texttt{Remember: zero is a valid score, meaning they are likely not feeling that emotion. You must score at least one emotion > 0.} \\
\\ 
\texttt{Your answer:} \\
            }
        }
    }
     \\ \\

    \theutterance \stepcounter{utterance}  
    & \multicolumn{4}{p{0.6\linewidth}}{
        \cellcolor[rgb]{0.561,0.737,0.561}{
            \makecell[{{p{\linewidth}}}]{
                \texttt{\tiny{[P1$\rangle$GM]}}
                \texttt{Relieved: 0} \\
\texttt{Indignant: 0} \\
\texttt{Guilty: 0} \\
\texttt{Bitter: 8} \\
            }
        }
    }
    & & \\ \\

    \theutterance \stepcounter{utterance}  
    & & & \multicolumn{2}{p{0.3\linewidth}}{
        \cellcolor[rgb]{0.9,0.9,0.9}{
            \makecell[{{p{\linewidth}}}]{
                \texttt{\tiny{[GM$|$GM]}}
                \texttt{Relieved: 0} \\
\texttt{Indignant: 0} \\
\texttt{Guilty: 0} \\
\texttt{Bitter: 8} \\
            }
        }
    }
    & & \\ \\

    \theutterance \stepcounter{utterance}  
    & & & \multicolumn{2}{p{0.3\linewidth}}{
        \cellcolor[rgb]{0.9,0.9,0.9}{
            \makecell[{{p{\linewidth}}}]{
                \texttt{\tiny{[GM$|$GM]}}
                \texttt{Relieved: 2} \\
\texttt{Indignant: 1} \\
\texttt{Guilty: 0} \\
\texttt{Bitter: 4} \\
            }
        }
    }
    & & \\ \\

    \theutterance \stepcounter{utterance}  
    & & & \multicolumn{2}{p{0.3\linewidth}}{
        \cellcolor[rgb]{0.9,0.9,0.9}{
            \makecell[{{p{\linewidth}}}]{
                \texttt{\tiny{[GM$|$GM]}}
                \texttt{game\_result = LOSE} \\
            }
        }
    }
    & & \\ \\

    \theutterance \stepcounter{utterance}  
    & & \multicolumn{4}{p{0.6\linewidth}}{
        \cellcolor[rgb]{0.95,0.95,0.95}{
            \makecell[{{p{\linewidth}}}]{
                \texttt{\tiny{[P1$\leftrightarrow$P]}}
                \texttt{Post{-}episode Reflexion (not sent to GM):} \\
\texttt{New plan: } \\
\\ 
\texttt{In my previous attempt, I failed to accurately predict Emma's emotional response. Upon re{-}reading the dialogue, I realized that I didn't consider the context of the conversation and the relationships involved. Emma is married to Steve, and Sean's confession of love puts her in a difficult position. } \\
\\ 
\texttt{To improve my prediction, I will focus on the following key points:} \\
\\ 
\texttt{1.  I will analyze the dialogue from Emma's perspective, considering her relationships and the potential consequences of her actions.} \\
\texttt{2.  I will pay attention to the emotional cues in the dialogue, such as Emma's hesitation and her attempts to clarify Sean's intentions.} \\
\texttt{3.  I will consider the power dynamics at play in the conversation, including Emma's marriage and Sean's confession of love.} \\
\\ 
\texttt{With these factors in mind, I will reassess the emotional response of Emma at the end of the dialogue.} \\
\\ 
\texttt{To do this, I will ask myself the following questions:} \\
\\ 
\texttt{{-} How does Emma's response to Sean's confession of love reflect her feelings about her marriage and her relationship with Steve?} \\
\texttt{{-} How does Emma's hesitation and her attempts to clarify Sean's intentions suggest that she is feeling?} \\
\texttt{{-} What are the potential consequences of Emma's actions in this situation, and how might she be feeling about those consequences?} \\
\\ 
\texttt{By considering these factors and asking myself these questions, I will be able to make a more accurate prediction of Emma's emotional response at the end of the dialogue.} \\
\\ 
\texttt{New plan: } \\
\\ 
\texttt{Relieved: 0} \\
\texttt{Indignant: 2} \\
\texttt{Guilty: 6} \\
\texttt{Bitter: 8} \\
            }
        }
    }
    & \\ \\

\end{supertabular}
}

\end{document}
